\Needspace{10\baselineskip}
\centerline{\textsc{Ceres.}}

\medskip
Monatliche Correspondenz. Vol. VIII, pp.\ 288—291 and pp.\ 369—371. September and October 1803.

\medskip
— — —

\medskip
Meanwhile, Dr.\ \textsc{Gauss} investigated what correction to his latest elements (VIII) followed from the Bremen observations and the three Palermo observations given on p.\ 94 of the July issue. A slight change was enough to bring these new observations into agreement with the old ones. Only the node had to be moved 3 minutes farther than for the observations made at Palermo in 1801. One may hope that a fuller calculation of the perturbations will eventually justify this. For now, here are these new elements (IX).

\begin{center}
\begingroup\fontsize{9}{11}\selectfont\setlength{\tabcolsep}{3pt}
\begin{tabular}{lr}
Epoch for Seeberg 1803 . . . . . . . . . . . . & $233^\circ\ 36'\ 3''1$ \\
Tropical motion in 365 days . . . . . . . . . & $78\ 8\ 7.2$ \\
Tropical motion in one day . . . . . . . . & $770''650$ \\
Logarithm of the semimajor axis . . . . . . . . & $0.4421516$ \\
Aphelion 1803 . . . . . . . . . . . . . . . & $326^\circ\ 33'\ 18''$ \\
Node 1803 . . . . . . . . . . . . . . . . . & $80\ 58\ 22$ \\
Eccentricity . . . . . . . . . . . . . . . . & $0.0788941$ \\
Inclination . . . . . . . . . . . . . . . . . . & $10^\circ\ 37'\ 54''$ \\
\end{tabular}\endgroup
\end{center}
\SourceLeaf{251}%

Here is how these elements agree with the three observations by \textsc{Piazzi} mentioned above:

\begin{center}
\begingroup\fontsize{9}{11}\selectfont\setlength{\tabcolsep}{3pt}\renewcommand{\arraystretch}{1.12}
\begin{tabular}{r|r|r|r|r|r}
 & & \multicolumn{2}{c|}{Calculated} & \multicolumn{2}{c}{Difference R—B} \\ \cline{3-6}
1803 & \shortstack{Mean time\\in Palermo} & \shortstack{Right ascension\\of $\CeresSymbol$} & \shortstack{Declination\\of $\CeresSymbol$} & in R.A. & in decl. \\ \hline
May 12 & $15^u\,53^m\,37^s7$ & $288^\circ\,19'\,17''9$ & $24^\circ\,36'\,22''0$ & $+\,2''9$ & $+\,9''0$ \\
--- 13 & $15\,49\,41.3$ & $288\,18\,57.3$ & $24\,39\,57.2$ & $-\,10.2$ & $-\,4.2$ \\
--- 14 & $15\,45\,41.9$ & $288\,18\,13.1$ & $24\,43\,36.7$ & $-\,1.9$ & $+\,5.6$ \\
\end{tabular}\endgroup
\end{center}

[The October issue contains the ephemeris calculated from \textsc{Gauss}’s elements IX. It runs from 30 April 1804 to 19 January 1805 at three-day intervals, for midnight at Seeberg. It gives the right ascension and declination to the nearest minute, the brightness to five decimal places, and also the mean time of transit across the meridian.]

\begin{flushright}
[\textsc{von Zach}.]
\end{flushright}

\bigskip
